\documentclass[10pt]{article}
\usepackage[margin=20mm]{geometry}
\usepackage{amsmath,amssymb,mathtools,fontspec,unicode-math,xurl,hyperref}
\IfFontExistsTF{DejaVu Serif}{\setmainfont{DejaVu Serif}}{\setmainfont{Latin Modern Roman}}
\setmathfont{latinmodern-math.otf}
\setlength{\parindent}{0pt}\setlength{\parskip}{6pt}\sloppy
\title{Technical Review of
Reflexive Calkin Algebras}
\author{AI-assisted technical review}
\date{}
\begin{document}
\maketitle
Abstract: The reviewed source is arXiv:2401.18037v1, dated 31 January 2024. Two local statements are false under their printed hypotheses. Finite-rank coordinate projections contradict the strictly singular quotient conclusion of Proposition 4.4, because its quantified subspace condition can be vacuous. A two-leaf tree contradicts Lemma 11.5(iii) for an arbitrary antichain. Sufficient repairs are proved here. The former proposition is not used in the central construction; the latter estimate is valid for the support restrictions actually used in Proposition 11.8. The demonstrated impact is Minor in each case. No counterexample to the central Calkin-algebra realization is established. The source passages supporting both findings were visually confirmed during recovery; the stated journal-version and long-proof limits remain.\par
\section*{Identity, scope, and terminology}
The source is P. Motakis and A. Pelczar-Barwacz, Reflexive Calkin algebras, arXiv:2401.18037v1 [MP]. The exact revision was submitted on 31 January 2024 at 18:01:20 UTC. The associated journal identifier is `10.4171/JEMS/1709`; the publisher records online publication on 24 September 2025 [EMS]. The acquired 85-page preprint has no verified equivalence to the accepted manuscript or journal version. Every statement below concerns this precise acquired version. Page numbers are one-based PDF pages, agreeing with printed pages. All 85 pages were read as extracted text.\par
Here \(\mathcal L(X)\) denotes bounded operators, \(\mathcal K(X)\) compact operators, and \(\mathcal{SS}(X)\) strictly singular operators. These ideals must not be interchanged. The main algebra is \(\mathcal L(X)/\mathcal K(X)\). The auxiliary Proposition 4.4 instead concerns the quotient by \(\mathcal{SS}(X)\). A projection may be nonzero and still belong to both ideals when its range is finite-dimensional.\par
The main result assumes a normalized unconditional basis of \(U\) with no \(c_0\) asymptotic version, in the alternating block-choice sense of Definition 1.1. It constructs a Banach space \(\mathfrak X_U\) and mutually annihilating infinite-rank projections \(I_s\) such that \[\mathcal L(\mathfrak X_U)
 =\mathcal K(\mathfrak X_U)\oplus[I_s:s\geq1]\oplus\mathbb CI,\] with \((I_s)\) equivalent to the basis of \(U\). The conclusion is a Banach-algebra realization of the unitization of \(U\) with coordinate multiplication. It does not assert a Hilbert-space \(C^*\)-algebra realization.\par
\section*{Checked proof chain and limits}
The review followed the following substantive dependencies. In Proposition 4.1, pp. 16–17, the lower estimate for the quotient by compact operators uses successive vectors annihilating earlier members of a norming set. The compact images admit close pairs; their difference gives the factor \(2C\) in the lower bound. No unjustified weak-null conclusion from a merely norming set is needed. The finite-dimensional cancellation in the bonding construction, particularly Proposition 6.12 and the estimates in 6.15–6.17, supplies the hypotheses of this proposition.\par
The Bourgain–Delbaen construction uses triangular extension functionals and self-determined subsets. The coordinate projections in Section 6 depend on the bounds for the extension operators, and the ground-coordinate cancellation controls the maximum of component norms. The lower and upper estimates then give the equivalence constant used in Theorem 10.1. Section 7 and Appendix A pass from absence of a \(c_0\) asymptotic version to a lower block \(\ell_q\) estimate and a dual upper \(\ell_p\) estimate. The quantifier order in the block choices is essential; an arbitrary lower estimate on single fixed blocks would not suffice.\par
The operator argument in Section 10 uses scalar-plus-compact behavior on each component, compactness between distinct components, and horizontal scalar-plus-compact behavior on the whole space. These are proved later by rapidly increasing sequences and dependent sequences in Sections 12–13. The final step, Theorem 10.8, is a norm-distance-to-compact argument for the tails of \(\sum_s\lambda_s I_s\); it does not assume bounded completeness of the basis of \(U\). The cancellation vector has norm at most \(8\), while the ground correction has norm at most \(4N_0+12\). The chosen tail bound \(\varepsilon/(8N_0+24)\) makes its image smaller than \(\varepsilon/2\), producing the required contradiction to horizontal compactness.\par
The tree-analysis definitions, weight separation, average estimates, and the two basic inequalities were read through their use in the operator chain. The local antichain defect below was traced to its actual use in Proposition 11.8. This review does not certify every constant in every rapidly increasing or dependent-sequence induction. The external Bourgain–Delbaen and asymptotic-structure theorems, and all omitted componentwise translations in Section 13, have not been independently reproved in full. These limits remain explicit; a complete textual reading is not a formal proof certification.\par
\section*{A vacuous hypothesis in the strictly singular quotient proposition}
Problem 1 (Finite-rank projections satisfy the printed hypotheses).  Proposition 4.4, p. 18, asserts an equivalence in \(\mathcal L(X)/\mathcal{SS}(X)\) for mutually annihilating projections \(I_s\), writing \(X_s=I_sX\). Its condition \((\gamma)\) quantifies over subspaces \(Y\) such that \(Y\cap X_s\) is infinite-dimensional for every \(s\). No infinite-rank hypothesis on the \(I_s\) is stated.\par
\subsection*{An exact counterexample}
Take \[U=X=\ell_1,\qquad u_s=e_s,\qquad I_sx=x_se_s,
 \qquad B=\beta=C=1.\] The basis of \(U\) is normalized and \(1\)-unconditional. The nonzero norm-one projections are mutually annihilating. For \(x_s=c_se_s\) and \(x_0\in\bigcap_{s=1}^N\ker I_s\), disjoint supports give \[\left\|\sum_{s=1}^Nx_s+x_0\right\|_1
 =\sum_{s=1}^N|c_s|+\|x_0\|_1
 \geq
 \left\|\sum_{s=1}^N\|x_s\|_1u_s\right\|_1.\] Thus condition \((\beta)\) holds with \(B=1\). Every \(X_s\) is one-dimensional. There is no subspace \(Y\) meeting the premise of \((\gamma)\), so its universal assertion is true, with the displayed positive \(\beta,C\).\par
Each \(I_s\) has finite rank and belongs to \(\mathcal{SS}(X)\). Consequently its quotient class is zero. A zero sequence cannot be equivalent to \((u_s^*)\), whose individual norms are one. Also \(0\neq I_1\in[I_s]\cap\mathcal{SS}(X)\), contradicting the claimed direct-sum conclusion. This checks every printed hypothesis; no numerical approximation or unverified structural example is involved.\par
The unquotiented sequence \((I_s)\) in \(\mathcal L(\ell_1)\) is indeed isometric to the unit-vector sequence spanning \(c_0\) in \(\ell_\infty\). The counterexample does not refute that separate part of the proposition.\par
\subsection*{A sufficient repair}
Add the requirement that every \(X_s\) be infinite-dimensional. Here is a direct verification of the quotient estimate under that additional premise and the original \((\beta),(\gamma)\).\par
For \(T\in\mathcal{SS}(X)\) and \(\delta>0\), choose a normalized basic sequence \((z_n)\) with basis constant at most \(2\), cycling through the spaces \(X_s\) so that each is visited infinitely often, and \[\|Tz_n\|\leq \delta\,2^{-n}/4.\] At each step the finitely many norming constraints used to continue a basic sequence leave an infinite-dimensional subspace of the chosen \(X_s\). Strict singularity guarantees a unit vector there with arbitrarily small image. The coefficient functionals have norms at most \(4\). Hence the closed span \(Y=[z_n]\) satisfies \(\|T|_Y\|\leq\delta\), and \(Y\cap X_s\) is infinite-dimensional for every \(s\).\par
Set \[P=\sum_{s=1}^N\lambda_sI_s,\qquad
 \Lambda=\left\|\sum_{s=1}^N\lambda_su_s^*\right\|_{U^*}.\] Condition \((\beta)\), applied to \(x_s=I_sx\) and \(x_0=x-\sum_{s=1}^NI_sx\), gives \[\|Px\|\leq\sum_s|\lambda_s|\|I_sx\|\leq B\Lambda\|x\|.\] Choose \(a_s\geq0\) with \(\|\sum_s a_su_s\|=1\) and \(\sum_s|\lambda_s|a_s=\Lambda\); finite-dimensional duality and unconditionality justify this choice. Apply \((\gamma)\) in \(Y\) to these \(a_s\) and \(b_s=\lambda_sa_s\). It supplies \(x=\sum_s a_sx_s-x_0\in Y\) with \[\|x\|\leq C,\qquad \|Px\|\geq\beta\Lambda.\] Since \(\|Tx\|\leq\delta\|x\|\), \[\beta\Lambda\leq C\bigl(\|P+T\|+\delta\bigr).\] Letting \(\delta\downarrow0\) and taking the infimum over \(T\) proves \[\frac{\beta}{C}\Lambda
 \leq \|P+\mathcal{SS}(X)\|
 \leq \|P\|\leq B\Lambda.\] These are the stated equivalence bounds. They extend to the closed spans. The identity is independent modulo \(\mathcal{SS}(X)\): outside the support of any finite sum of projections, its restriction to \(X_s\) is the identity on an infinite-dimensional space. Approximation by finite sums rules out membership in the closed span modulo \(\mathcal{SS}(X)\). Thus the added infinite-dimensionality premise also gives the directness asserted for the identity extension.\par
\subsection*{Impact}
The source explicitly says that Proposition 4.4 is not used. The actual construction proceeds through Proposition 4.1, which has different hypotheses and the compact-operator quotient. This defect is therefore Minor at the demonstrated scope. The strengthened proposition above is verified, but neither author adoption nor persistence in the journal text is asserted.\par
\section*{An antichain need not capture the input}
Problem 2 (Lemma 11.5(iii) lacks a support condition).  On p. 53, for an arbitrary antichain \(\mathcal B\) in a tree analysis \((f_t)_{t\in\mathcal T}\), the source asserts \[|f(x)|\leq
 \left(\sum_{t\in\mathcal B}|f_t(x)|^q\right)^{1/q}
 \quad\text{for every }x\in c_{00}.\] An antichain can omit an entire active branch.\par
\subsection*{An admissible two-leaf tree}
Use \(p=q=2\), \(m_1=16\), \(n_1=256\), and the norming set \(W_q[(\mathcal A_{4n_j},1/m_j)]\). Take \(U=\ell_2\) with its coordinate blocking, so Assumption 7.3 holds with \(\theta_U=1\) and \(p=q=2\). These values satisfy the initial growth requirements in Assumption 8.3; subsequent values can be chosen recursively to meet all the stated lower growth bounds. The first operation permits the two-coordinate functional \[f=\frac{e_1^*+e_2^*}{m_1(4n_1)^{1/2}}
   =\frac{e_1^*+e_2^*}{512}.\] Its tree has a root and two terminal children. Let \(\mathcal B\) consist only of the child labeled \(e_1^*\), and let \(x=e_2\). This is a nonempty antichain, yet \[|f(x)|=\frac1{512}>0,\qquad
 \left(\sum_{t\in\mathcal B}|f_t(x)|^2\right)^{1/2}=0.\] An independently executed rational calculation reproduces these exact values. The mathematical counterexample is the explicit permitted tree, not an extrapolation from that calculation.\par
\subsection*{Correct support-sensitive estimates}
Write \[S_{\mathcal B}=\bigcup_{t\in\mathcal B}\operatorname{supp}f_t.\] The claimed estimate is valid if \(\operatorname{supp}f\cap\operatorname{supp}x\subseteq S_{\mathcal B}\). Equivalently, every root-to-terminal path whose leaf acts nontrivially on \(x\) must meet \(\mathcal B\).\par
Equation reference: restricted. Indeed, a nonterminal node with \(d\leq n\) children, operation weight \(0<\theta\leq1\), and factor \(n^{-1/p}\) satisfies \[|f_t(x)|
 \leq\theta n^{-1/p}\sum_r|f_r(x)|
 \leq\theta(d/n)^{1/p}
      \left(\sum_r|f_r(x)|^q\right)^{1/q}
 \leq\left(\sum_r|f_r(x)|^q\right)^{1/q}.\] Stop this induction at the antichain. Every other terminal contribution is zero by the support condition. The same argument includes \(p=\infty,q=1\) by the triangle inequality. In particular, for every antichain and every \(x\), \[
 \left|f(x|_{S_{\mathcal B}})\right|
 \leq\left(\sum_{t\in\mathcal B}|f_t(x)|^q\right)^{1/q}.\] If \(K=\operatorname{supp}f\setminus S_{\mathcal B}\), the unrestricted version is \[|f(x)|\leq |(f|_K)(x)|
       +\left(\sum_{t\in\mathcal B}|f_t(x)|^q\right)^{1/q}.\] The missing term cannot generally be discarded.\par
\subsection*{The actual use in Proposition 11.8}
The proof on pp. 56–57 separates precisely the missing set \(K\). Its blocks have \(\operatorname{supp}y_i\subseteq I_i\), with the \(I_i\) successive and disjoint. The antichain is partitioned into \(\mathcal B_i\), and the source explicitly defines \[J=\bigcup_i\bigcup_{t\in\mathcal B_i}(\operatorname{supp}f_t\cap I_i),
 \qquad f'=f|_J.\] For \[D_i=\bigcup_{\ell\ne i}\bigcup_{t\in\mathcal B_\ell}\operatorname{supp}f_t,
 \qquad z_i=y_i|_{D_i},\] disjointness of the antichain supports gives the exact identity \[(f-f')(y_i)=(f|_K)(y_i)+f(z_i).\] Applying [restricted] (\#restricted) to the cross antichain yields \[|f(z_i)|\leq
 \left(\sum_{t\in\bigcup_{\ell\ne i}\mathcal B_\ell}
                    |f_t(y_i)|^q\right)^{1/q}.\] Thus the local separation estimate used there is valid with constant one and no changed assumptions on Proposition 11.8. Its separately bounded \(K\) contribution is retained. This verifies the precise use affected by the overbroad lemma; it is not a fresh certification of every subsequent induction.\par
The lemma as printed is false, but the needed restricted estimate and its actual application are established above. The demonstrated defect is consequently Minor. A different later use without support coverage would need a separate check; no such uncorrected use has been demonstrated here.\par
\section*{Assessment and remaining dependencies}
Location | Assessment | Verified repair and impact\par
Proposition 4.4 | Minor | Add infinite-dimensional ranges; quotient estimates proved. The main construction does not use this proposition.\par
Lemma 11.5(iii) | Minor | Retain uncovered support or require coverage. The actual restricted use in Proposition 11.8 survives with the same constant.\par
Main realization | No defect established in checked steps | Long dependent-sequence estimates and external structure results are not fully independently certified.\par
Scientific decision remains pending, and no human review is recorded.\par
The acquisition record and a bounded follow-up search of the arXiv record, publisher page, and author publication page did not identify a separate supplement, paper-specific implementation, or correction. This is not proof of nonexistence. The publisher page requires a subscription; no access control was bypassed. No author code was run. Absence of code is not a defect in this mathematical paper.\par
9 P. Motakis and A. Pelczar-Barwacz, Reflexive Calkin algebras, arXiv:2401.18037v1, 31 January 2024. https://arxiv.org/abs/2401.18037v1. EMS Press, metadata for DOI 10.4171/JEMS/1709, consulted 27 October 2026. https://ems.press/journals/jems/articles/14299156. P. Motakis, publication list, consulted 27 October 2026. https://pmotakis.mathstats.yorku.ca/publications/.\par
\section*{Supporting checks}
752d5908f677c0e56cffdf372e63f7061433ad0d39fdd2b1da1129b74323ca63.\par
The original initial executor read the complete extracted source, checked the final arguments, wrote this report, and performed the recorded computations and compilation. It did not acquire files, inspect images, run software, write the report, or conduct independent acceptance. The original executor checked the source’s explicit definition \(f'=f|_J\) before using the local application repair.\par
The controller’s UI observation and this runtime’s self-description are distinct records; no backend-model equivalence is claimed. The original image-reader EMFILE failures are preserved. Original rational-calculation outputs were reused with unchanged hashes, not rerun or reattributed. The recovered report clarifies the admissible \(U=\ell_2\) parameter choice. Exact current compilation, dependency, and report-page inspection facts are retained separately. This continuation is initial-review recovery, not independent acceptance.\par
\end{document}
